\documentclass[a4paper,11pt]{article}
\usepackage{exam-style-842-unit}
\examinehead{842 信息化概论 · 第 2 单元练习册}{练习一 参考答案与解析}
\begin{document}

\unitanswerscover{842 信息化概论 · 单元练习册}{第 2 单元参考答案}{练习一（数据管理与数据库基础）的答案与解析；每道判断题均给出改正后的正确表述。}

\answerpart{一、单项选择题（10 题，每题 3 分，共 30 分）}

\q B。\begin{analysis}数据管理的发展一般划分为四个阶段：人工管理（20 世纪 50 年代中期以前）$\rightarrow$ 文件系统（20 世纪 50 年代末至 60 年代中期）$\rightarrow$ 数据库系统（20 世纪 60 年代末以后）$\rightarrow$ 高级数据库技术（20 世纪 80 年代以后）。（大纲「二、计算机与软件技术」第 2 条；教材 3.2.2 第 34—35 页指出数据库系统萌芽于 20 世纪 60 年代、此前由文件系统管理数据。）A 把人工管理排在最后；C 把数据库系统提到文件系统之前；D 把高级数据库技术插在数据库系统之前，均与阶段先后矛盾。\end{analysis}

\q B。\begin{analysis}这是教材 3.2.2（第 34 页）对数据库的定义：数据库是以一定方式储存在一起、能为多个用户共享、具有尽可能小的冗余度、与应用程序彼此独立的数据集合。A 数据是描述事物属性的符号记录；C 数据库管理系统是管理数据的系统软件；D 数据库系统是由数据库、DBMS、应用程序和 DBA 构成的整体，范围都比数据库大。\end{analysis}

\q C。\begin{analysis}教材 3.2.2（第 34 页）原文：数据库系统的萌芽出现于 20 世纪 60 年代，当时计算机开始广泛地应用于数据管理，对数据的共享提出了越来越高的要求，传统的文件系统已不能满足需要，DBMS 应运而生。A、B 两个时间计算机尚处于科学计算阶段；D 已是关系数据库普及之后。\end{analysis}

\q A。\begin{analysis}层次模型是一种树结构模型，把数据按自然的层次关系组织起来，以反映数据之间的隶属关系，结点代表数据记录、连线描述从属关系，只能表示一对多联系（教材 3.2.2，第 35 页）。B 网状模型用有向图结构、可表示多对多联系；C 关系模型用二维表表示实体及其联系；D 概念模型是面向用户的 E-R 模型，不是数据库的逻辑数据模型。\end{analysis}

\q C。\begin{analysis}关系模型可以简单、灵活地表示各种实体及其关系，在关系数据库系统中对数据的操作通过关系代数实现，具有严格的数学基础，因而 20 世纪 80 年代以后推出的 DBMS 几乎都是关系型的（教材 3.2.2，第 35 页）。A 层次模型、B 网状模型都用指针或路径导航实现联系，缺乏严格的集合论基础；D 面向对象模型不在教材所列的三种传统数据模型之内。\end{analysis}

\q A。\begin{analysis}SQL 是结构化查询语言（Structured Query Language）的缩写，是一种特殊目的的编程语言，用于存取数据以及查询、更新和管理关系数据库系统（教材 3.2.4，第 37—38 页）。B、C、D 都是对 SQL 全称的臆造；注意 SQL 是非过程化语言，用户不必指定数据的存放方法。\end{analysis}

\q A。\begin{analysis}结构化数据可以组织成行列结构，是可识别、可精确定位的数据，其两个重要特征是排序和索引：排序使数据有结构、能快速定位，索引是更特殊的排序，能更快定位数据（教材 3.2.5，第 38 页）。B 加密压缩、C 共享传输、D 备份恢复都是数据处理与保护手段，不是结构化数据区别于非结构化数据的结构特征。\end{analysis}

\q B。\begin{analysis}视图是从一个或多个基本表（或其他视图）导出的表，它是一个虚表，数据库中只存放视图的定义，不存放视图对应的数据，视图中的数据随基本表数据的变化而变化（大纲「数据库原理与基本概念」；教材 3.2 未展开视图，说明见本单元 README）。A 基本表是实际存储数据的表；C 索引是加快检索的存储结构；D 数据字典存放的是数据库的定义信息，均不是视图。\end{analysis}

\q C。\begin{analysis}数据库的三级模式为外模式（用户模式、子模式）、概念模式（模式）和内模式（存储模式）：外模式描述用户能看到的数据视图，概念模式描述全体数据的逻辑结构，内模式描述数据的物理存储结构和存取方法（大纲补充；通行数据库理论）。A 外模式面向用户；B 概念模式面向数据库整体；D 子模式是外模式的别称。\end{analysis}

\q D。\begin{analysis}事务具有原子性、一致性、隔离性和持久性：原子性指事务的操作要么全做要么全不做；一致性指事务把数据库从一个一致状态带到另一个一致状态；隔离性指并发事务互不干扰；持久性指事务一旦提交，对数据库的修改就是永久的，即使系统故障也不会丢失（大纲补充；通行数据库理论）。A、B、C 三项分别对应另外三个特性，与题干描述不符。\end{analysis}

\answerpart{二、多项选择题（4 题，每题 4 分，共 16 分）}

\q ABC。\begin{analysis}A、B、C 三项正是教材 3.2.2（第 34 页）数据库定义中的三个要点：能为多个用户共享、具有尽可能小的冗余度、与应用程序彼此独立。D 错误：数据库中的数据由 DBMS 统一组织和管理，不是以彼此独立的文件形式零散存放；把数据库等同于“一堆文件”正是文件系统阶段的特征。\end{analysis}

\q ABC。\begin{analysis}数据模型是对现实世界数据特征的抽象，其三要素是数据结构、数据操作和完整性约束（大纲「数据库原理与基本概念」；通行数据库理论）。A 数据结构描述数据库的组成对象以及对象之间的联系；B 数据操作指对数据库中各种对象允许执行的操作集合；C 完整性约束是一组完整性规则。D 错误：存储结构属于内模式（物理模型）的内容，不是数据模型的三要素。\end{analysis}

\q ABC。\begin{analysis}A 正确：外模式是用户能看到并使用的局部数据视图；B 正确：概念模式描述全体数据的逻辑结构；C 正确：模式/内模式映像保证物理独立性，内模式改变时只需修改映像；D 错误：外模式/模式映像保证的是逻辑独立性，即概念模式改变时应用程序可以不变。\end{analysis}

\q ABCD。\begin{analysis}事务的 ACID 特性即原子性（Atomicity）、一致性（Consistency）、隔离性（Isolation）和持久性（Durability），四项全对。（大纲「数据库原理与基本概念」；通行数据库理论）\end{analysis}

\answerpart{三、判断题（6 题，每题 2 分，共 12 分）}

\q ×。改正：数据库应具有尽可能小的冗余度，冗余度过大会造成数据重复存储、浪费空间并引起更新异常。\begin{analysis}教材 3.2.2（第 34 页）的数据库定义明确要求“具有尽可能小的冗余度”，并要求数据与应用程序彼此独立。合理控制冗余是数据库优于文件系统的重要方面。\end{analysis}

\q ×。改正：关系模型中主键的取值必须唯一，且不能取空值，这正是实体完整性的要求。\begin{analysis}主键用于唯一标识关系中的每一个元组，若允许为空或重复，就无法区分不同元组。若主键由多个属性组成（组合键），则每个属性都不能为空。\end{analysis}

\q √。\begin{analysis}视图是从一个或多个基本表（或其他视图）导出的虚表，数据库中只存放视图的定义，不存放视图对应的数据；视图中的数据随基本表数据的变化而变化，对视图的更新最终转换为对基本表的更新。（大纲补充；通行数据库理论）\end{analysis}

\q ×。改正：外模式/模式映像保证数据的逻辑独立性，模式/内模式映像保证数据的物理独立性。\begin{analysis}两级映像与两级独立性一一对应：概念模式改变时，只需修改外模式/模式映像，外模式和应用程序不变，这是逻辑独立性；内模式改变时，只需修改模式/内模式映像，概念模式和应用程序不变，这是物理独立性。\end{analysis}

\q √。\begin{analysis}数据库系统阶段的数据具有整体结构化特征：数据不再只面向某一应用，而是面向整个系统组织，由 DBMS 统一管理，具有共享性高、冗余度小、易扩充和独立性强等特点（教材 3.2.2，第 34—35 页）。\end{analysis}

\q √。\begin{analysis}教材 3.2.2（第 35 页）指出：在关系数据库系统中，对数据的操作是通过关系代数实现的，具有严格的数学基础。这也是关系模型发展迅速、成为主流的重要原因。\end{analysis}

\answerpart{四、配对题（6 对，每对 1 分，共 6 分）}

\q （1）—B\quad （2）—C\quad （3）—D\quad （4）—A\quad （5）—F\quad （6）—E。
\begin{analysis}层次模型是树形结构、表示一对多联系；网状模型是有向图结构、结点间联系任意、可表示多对多联系；关系模型用二维表并以关系代数操作；视图是从基本表或视图导出的虚表；主键唯一标识元组且非空；外键是参照另一关系主键的属性组。（数据模型部分见教材 3.2.2，第 35 页；视图与键见大纲补充）\end{analysis}

\answerpart{五、简答题（3 题，每题 8 分，共 24 分）}

\q 参考答案要点：
（1）人工管理阶段（20 世纪 50 年代中期以前）：计算机主要用于科学计算，没有专门的数据管理软件，数据由程序员以手工方式管理；数据不长期保存、不共享，冗余度大，程序与数据之间没有独立性。
（2）文件系统阶段（20 世纪 50 年代末至 60 年代中期）：操作系统中的文件系统把数据组织成文件长期保存在外存上，程序与数据之间有一定独立性；但数据仍面向具体应用，共享性差、冗余度大、独立性差，数据联系靠程序实现。
（3）数据库系统阶段（20 世纪 60 年代末以后）：出现数据库管理系统，数据整体结构化并面向整个系统组织，共享性高、冗余度小；DBMS 统一提供安全性、完整性、并发控制和恢复等数据控制功能；三级模式两级映像保证数据的逻辑独立性和物理独立性。（教材 3.2.2，第 34—35 页）
（4）高级数据库技术阶段（20 世纪 80 年代以后）：为适应应用规模扩大和数据形式多样化，分布式数据库、面向对象数据库、多媒体数据库、数据仓库和 NoSQL 等新技术相继出现，数据管理由集中走向分布、由结构化数据扩展到非结构化数据。
\begin{analysis}得分点：四个阶段名称及大致年代 4 分，各阶段主要特点 4 分。常见失分：只写“文件系统—数据库系统”两个阶段；把“数据整体结构化、共享性高”错记到文件系统阶段；漏答高级数据库技术阶段。前三个阶段以教材 3.2.2 为准，第四阶段属大纲要求、教材未展开。\end{analysis}

\q 参考答案要点：
（1）外模式（又称用户模式、子模式）：面向具体用户和应用，是用户能够看到并使用的那部分数据的逻辑结构，一个数据库可以有多个外模式。
（2）概念模式（简称模式）：面向数据库整体，描述全体数据的逻辑结构和特征，是全部数据的公共视图，独立于具体的存储结构。
（3）内模式（又称存储模式）：面向存储设备，描述数据的物理存储结构和存取方法，如记录如何存储、是否建立索引、数据如何压缩等。
（4）外模式/模式映像保证逻辑独立性：当概念模式改变时，数据库管理员只需修改外模式/模式映像，使外模式保持不变，从而应用程序不必修改。
（5）模式/内模式映像保证物理独立性：当内模式的存储结构改变时，只需修改模式/内模式映像，使概念模式保持不变，从而应用程序不必修改。
\begin{analysis}得分点：三级模式各 1 分共 3 分，两级映像及对应独立性各 2 分共 4 分，总结数据独立性的意义 1 分。常见失分：把两种映像对应的独立性写反；只背“三级模式”名称而不说明各自面向的对象。教材 3.2 节未展开三级模式，本答案依据大纲“数据库原理与基本概念”及通行数据库理论（王珊《数据库系统概论》口径）。\end{analysis}

\q 参考答案要点：
（1）视图是从一个或多个基本表（或其他视图）导出的表，它是一个虚表：数据库中只存放视图的定义，不存放视图对应的数据，视图中的数据随基本表数据的变化而变化。
（2）视图的优点：
（1）简化操作：把复杂的多表连接和统计条件定义为一个视图，用户只需对视图作简单查询；
（2）提供逻辑独立性：当基本表结构发生变化时，可通过修改视图定义使应用程序不必改动；
（3）实现安全保护：只把视图中允许用户看到的行和列暴露给用户，隐藏机密数据；
（4）便于数据共享与定制：同一基本表可以按不同用户的需要定义多种视图，从不同角度提供数据；
（5）视图机制使不同用户只能访问其权限范围内的数据，简化了权限管理。
（3）注意：对视图的更新最终要转换成对基本表的更新，含聚集函数、分组、多表连接等的视图通常不可更新。
\begin{analysis}得分点：视图定义（虚表、只存定义）3 分，写出四点以上优点 4 分，更新限制 1 分。常见失分：说视图在数据库中物理保存了一份数据；只答“方便查询”而漏掉安全性、逻辑独立性和数据共享。教材 3.2 节只讲数据库与 SQL，未展开视图，本答案依大纲“数据库原理与基本概念”及通行数据库理论作答（第 02 单元 README 中已注明）。\end{analysis}

\answerpart{六、综合题（1 题，共 12 分）}

\q 参考答案要点：
（1）E-R 模型（3 分）：实体有三个——种植户（属性：户编号、姓名、村组）、地块（属性：地块编号、面积、土壤类型）、作物（属性：作物编号、名称、生育期）。联系有两个：种植户与地块之间为“承包”联系，一个种植户可以承包多个地块、一个地块只属于一个种植户，属于一对多（1:n）联系；地块与作物之间为“种植”联系，一个地块可以种植多种作物、一种作物也可种在多个地块上，属于多对多（m:n）联系，可在联系上附加“种植面积”“季节”等属性。E-R 图中实体用矩形、属性用椭圆、联系用菱形表示。
（2）关系模式（4 分）：
种植户（户编号，姓名，村组），主键为户编号；
地块（地块编号，面积，土壤类型，户编号），主键为地块编号，外键为户编号（参照种植户的户编号）；
作物（作物编号，名称，生育期），主键为作物编号；
种植（地块编号，作物编号，种植面积，季节），主键为（地块编号，作物编号），外键为地块编号（参照地块）、作物编号（参照作物）。
（3）三类完整性（3 分）：
实体完整性：各关系的主键必须唯一且不能为空，如户编号、地块编号、作物编号不能为空，种植关系的组合主键（地块编号，作物编号）也不能为空；
参照完整性：地块关系中的户编号必须是种植户关系中已存在的户编号；种植关系中的地块编号、作物编号必须分别在地块关系和作物关系中存在，不得引用不存在的实体；
用户定义完整性：面积必须大于 0，土壤类型只能取规定的分类值，生育期天数应在合理范围内，种植面积不得超过地块面积等。
（4）视图的作用（2 分）：可为不同用户建立只含其所需数据的视图，如为乡镇管理部门建立“各村种植面积汇总视图”、为合作社建立“本社地块与作物视图”，从而简化查询、隐藏户主联系方式等敏感字段实现安全保护，并在基本表结构变化时通过修改视图定义保持应用程序稳定，即提供逻辑独立性。
\begin{analysis}得分点：E-R 实体、属性、联系类型 3 分；关系模式及主键外键 4 分；三类完整性 3 分；视图作用 2 分。常见失分：把“承包”联系误判为多对多；m:n 联系转换时漏建联系表，或联系表漏写组合主键；参照完整性只写“外键”而不说明参照方向。E-R 与关系模式转换属大纲“数据库原理与基本概念”要求，教材 3.2 节未展开，已在本单元 README 注明。\end{analysis}

\end{document}
